\section{Byzantine integrity: a geometric secure-identification bound}
\label{sec:byzantine}

An autonomous constellation cannot assume its peers are honest. A compromised or
malfunctioning node may inject \emph{arbitrary}, possibly colluding, biases on all the
measurements it participates in. This section asks how many such Byzantine peers the network
can detect and identify from its internal geometry alone. The answer is a
\emph{secure-identification} bound: a block-structured instance of the sparse-observability /
error-correcting-code bound of secure state estimation \citep{fawzi2014,shoukry2016,%
candes2005,donoho2003}. The bound is \Cited{}; its instantiation on the coupled-datum lunar
$\Gmat$ is \Proven{}; the tolerance \emph{counts} are \Validated{} pipeline outputs, and any
reading of them as an operational flight guarantee is \Modelled{}.

\subsection{Peer signatures and the effective (parity-projected) block}
\label{sec:peersig}

A peer fault at node $j$ has signature $\Bmat_j\in\R^{84\times m_j}$, the column-block
selecting node $j$'s $m_j$ incident measurement rows (Section~\ref{sec:taxonomy}); the
injected bias is $\bvec=\Bmat_j\bm v$ for an arbitrary $\bm v\in\R^{m_j}$. A coalition
$T\subseteq\{1,\dots,M\}$ of Byzantine peers injects $\bvec=\Bmat_T\bm v_T$, where
$\Bmat_T=[\Bmat_j]_{j\in T}$ stacks the coalition's column-blocks. What the parity test sees
is the \emph{effective signature}
\begin{equation}
  \Gbar_j \;:=\; \Pperp\Bmat_j,
  \qquad
  \Gbar_T \;=\; \Pperp\Bmat_T ,
  \label{eq:effsig}
\end{equation}
the peer block after the parity projector \eqref{eq:pperp}. A coalition is invisible exactly
when its injected bias can be hidden in $\rng(\Gmat)$, i.e.\ when $\Gbar_T$ is column-rank
deficient.

\begin{definition}[Block-spark]
\label{def:blockspark}
The \emph{block-spark} of the family $\{\Gbar_j\}_{j=1}^{M}$ is the smallest number of
peer-blocks whose union is column-rank deficient:
\begin{equation}
  \spark_{\mathrm b}\bigl(\{\Gbar_j\}\bigr)
  \;=\;\min\Bigl\{\,|T|\;:\;\Gbar_T=\Pperp\Bmat_T\ \text{is column-rank deficient}\,\Bigr\}.
  \label{eq:blockspark}
\end{equation}
It is the block analogue of the spark of a dictionary \citep{donoho2003} (the smallest number
of columns that are linearly dependent), lifted from single columns to peer column-blocks.
\end{definition}

\subsection{The detect / identify bound}
\label{sec:bound}

\begin{theorem}[Geometric secure detect / identify bound]
\label{thm:byzantine}
Let $\{\Gbar_j\}$ be the effective peer signatures \eqref{eq:effsig} on the free-network
$\Gmat$. Then:
\begin{enumerate}
  \item[(i)] \textbf{Detection.} All coalitions of size $\le f$ are detectable if and only if
  $\Gbar_T$ is full column rank for every $|T|\le f$, equivalently
  $\spark_{\mathrm b}(\{\Gbar_j\})>f$, equivalently $\operatorname{span}(\Bmat_T)\cap
  \rng(\Gmat)=\{\bm0\}$ for every $|T|\le f$.
  \item[(ii)] \textbf{Identification.} All coalitions of size $\le f$ are uniquely
  identifiable if and only if $\Gbar_U$ is full column rank for every $|U|\le 2f$,
  equivalently $\spark_{\mathrm b}(\{\Gbar_j\})>2f$ (a coalition minimum distance $\ge2f+1$).
\end{enumerate}
\end{theorem}

\begin{proof}
\emph{(i).} A size-$\le f$ coalition $T$ is undetectable iff there is a nonzero injection
$\Bmat_T\bm v$ with $\Pperp\Bmat_T\bm v=\bm0$, i.e.\ iff $\Gbar_T=\Pperp\Bmat_T$ has a
nontrivial kernel, i.e.\ iff $\Gbar_T$ is column-rank deficient. So \emph{all}
$|T|\le f$ are detectable iff every such $\Gbar_T$ is full column rank, which by
Definition~\ref{def:blockspark} is $\spark_{\mathrm b}>f$. By
Proposition~\ref{prop:projector}, $\Pperp\Bmat_T\bm v=\bm0\iff\Bmat_T\bm v\in\ker\Pperp=
\rng(\Gmat)$, giving the intersection form. \emph{(ii).} Two coalitions $T_1,T_2$ with
$|T_i|\le f$ produce the same residual iff $\Pperp(\Bmat_{T_1}\bm v_1-\Bmat_{T_2}\bm v_2)=
\bm0$, i.e.\ iff some injection on $T_1\cup T_2$ ($|T_1\cup T_2|\le2f$) is undetectable.
Unique identification of every $|T|\le f$ therefore requires every $|U|\le2f$ union to be
detectable, i.e.\ $\Gbar_U$ full column rank for all $|U|\le2f$, i.e.\ $\spark_{\mathrm
b}>2f$. This is the block form of the classical $2f$-sparse-observability / minimum-distance
$\ge2f+1$ condition of \citep{fawzi2014,shoukry2016,candes2005}.
\end{proof}

\paragraph{The detect condition is full column rank, not ``$\Pperp\Bmat_T\neq\bm0$.''}
It is tempting to declare a coalition detectable as soon as its projected block is nonzero.
This is \emph{wrong}. A block $\Gbar_T=\Pperp\Bmat_T$ can be nonzero yet column-rank
deficient, and the adversary then hides a fault in the null combination---an injection
$\Bmat_T\bm v\in\rng(\Gmat)$ that produces \emph{zero} residual even though $\Gbar_T\neq
\bm0$ on other injections. The correct detect condition is that every $\le f$-block union be
\emph{full column rank} (span disjoint from $\rng(\Gmat)$); the identify condition raises
this to $2f$.

\subsection{The Lamport distinction: this is not a \texorpdfstring{$3f+1$}{3f+1} consensus bound}
\label{sec:lamport}

We adopt the Byzantine fault \emph{model}---arbitrary, colluding, worst-case node
behaviour---from Lamport--Shostak--Pease \citep{lamport1982}, and \emph{only} the model. The
bound of Theorem~\ref{thm:byzantine} is a \emph{measurement-redundancy / coding} bound: it
counts how many arbitrarily corrupted \emph{measurement blocks} a geometric estimator can
detect ($\spark_{\mathrm b}>f$, i.e.\ $f+1$ block-redundancy, for detection) and correct
($\spark_{\mathrm b}>2f$ for identification), exactly as in secure state estimation and
error-correcting codes
\citep{fawzi2014,shoukry2016,candes2005,donoho2003}. It is \emph{not}, and must not be read
as, the $3f+1$ threshold for Byzantine agreement in a message-passing consensus network
\citep{lamport1982}. Those are different problem classes: consensus concerns how many
faulty \emph{processes} exchanging \emph{messages} still permit agreement; ours concerns how
many faulty \emph{measurement blocks} in a \emph{static estimation} geometry still permit
fault correction. The two thresholds ($3f+1$ versus $\spark_{\mathrm b}>2f$) are not
comparable, and we draw no analogy between them.

\subsection{Redundancy raises tolerance; broad coverage destroys it}
\label{sec:counts}

The bound of Theorem~\ref{thm:byzantine} is governed by the network's block-spark, and the
pipeline computes it by exact enumeration for the representative geometry (the counts are
\Validated{}; their reading as an operational guarantee is \Modelled{}).

\paragraph{Redundant network.}
A network of six independent single-measurement peers---each owning one linearly independent
row---has $\spark_{\mathrm b}=7$, giving $f_{\mathrm{detect}}=6$ and $f_{\mathrm{identify}}=3$
(the largest $f$ with $\spark_{\mathrm b}>2f$). Redundancy buys tolerance: adding independent
cross-links or providers raises the block-spark and hence $f$.

\paragraph{Coverage fragility and the deficiency-of-five self-blindness.}
The opposite regime is instructive and is the paper's cautionary result. A \emph{broad}-%
coverage peer that owns \emph{all} of node~2's incident rows (all $21$ of them) has
$\spark_{\mathrm b}=1$ and $f_{\mathrm{detect}}=0$: a \emph{single} such peer is already
\textbf{undetectable}. A \emph{narrow} peer owning one measurement has $\spark_{\mathrm b}=2$
and $f_{\mathrm{detect}}=1$, and is detectable. The mechanism is precise: a peer owning all
of a node's measurements can inject a \emph{uniform self-bias} equal to a shift of that
node's own five local degrees of freedom---its three position components, its clock offset,
and its clock rate. That injection is of the form $\Gmat(\text{shift of node's local
block})\in\rng(\Gmat)$, hence undetectable by Theorem~\ref{thm:undetectable}: \emph{a node is
blind to a fault in its own five local DOF}. Broad measurement coverage by a single peer is
therefore the fragility, not the strength. This is why the detect condition is full column
rank of the union and not mere nonzeroness: the broad-coverage block is nonzero after
projection on most injections, yet column-rank deficient on the self-bias direction.

\begin{table}[t]
\centering
\caption{Byzantine block-spark on the representative network (real DE440 orientation).
Counts are \Validated{} (matched exactly by the independent oracle, Section~\ref{sec:anchor});
their reading as an operational flight tolerance is \Modelled{}. $f_{\mathrm{detect}}=
\spark_{\mathrm b}-1$; $f_{\mathrm{identify}}=\lfloor(\spark_{\mathrm b}-1)/2\rfloor$.}
\label{tab:byzantine}
\small
\setlength{\tabcolsep}{7pt}
\begin{tabular}{@{}lccc@{}}
\toprule
\textbf{Peer structure} & $\spark_{\mathrm b}$ & $f_{\mathrm{detect}}$ & $f_{\mathrm{identify}}$ \\
\midrule
Six independent single-measurement peers & $7$ & $6$ & $3$ \\
Narrow peer (one incident measurement)   & $2$ & $1$ & $0$ \\
Broad peer (all $21$ rows of node 2)     & $1$ & $0$ & $0$ \\
\bottomrule
\end{tabular}
\end{table}

\subsection{Provider faults: the common / differential split}
\label{sec:providersplit}

The most structured Byzantine fault is a \emph{provider}-wide bias---a coordinated error
across all nodes of one provider (Section~\ref{sec:taxonomy}). Such a fault splits into a
part all providers share and a part that distinguishes them.

\begin{proposition}[Provider common / differential detectability]
\label{prop:provider}
A provider-wide bias decomposes as $\bvec=\bvec_{\mathrm{common}}+\bvec_{\mathrm{diff}}$,
where the common part---shared identically by every provider---lies in $\rng(\Gmat)$ and is
therefore \emph{internally undetectable} (it is the common-mode class of
Remark~\ref{rem:commonmode}, requiring an external tie, and is the common component of the
inter-ephemeris floor of \citep{baweja2026interop}); and the differential part---provider $A$
disagreeing with provider $B$---has nonzero parity, $\Pperp\bvec_{\mathrm{diff}}\neq\bm0$,
and is \emph{internally detectable}. Consequently: with two providers the network
\emph{detects} disagreement but cannot \emph{identify} the culprit (the two-clock problem);
with three or more providers it identifies the culprit via block-spark $>2f$
(Theorem~\ref{thm:byzantine}).
\end{proposition}

\begin{proof}
The common part shifts every provider's realization by the same amount---a whole-network
datum move---so it is $\Gmat(\cdot)\in\rng(\Gmat)$ and $\Pperp\bvec_{\mathrm{common}}=\bm0$
by Theorem~\ref{thm:undetectable}. The differential part is supported on a proper subset of
provider blocks and is not a whole-network datum move, so it does not lie in $\rng(\Gmat)$
and has nonzero residual. Detection versus identification then follows from
Theorem~\ref{thm:byzantine}: two providers give two blocks whose disagreement is one
detectable direction but with no third block to break the tie ($\spark_{\mathrm b}=2$,
$f_{\mathrm{identify}}=0$); three or more give $\spark_{\mathrm b}>2$, enough to identify a
single faulty provider.
\end{proof}

The pipeline reports the number of internally detectable cross-provider rows as a function
of provider count (\Modelled{}): $1\to0$, $2\to48$, $4\to72$ (Table~\ref{tab:provider}). One
provider offers nothing to cross-check; two expose disagreement; four expose more. The count
grows sublinearly because the common part---the shared datum move---is never detectable no
matter how many providers share it.

\begin{table}[t]
\centering
\caption{Internally detectable cross-provider rows versus provider count (\Modelled{}). One
provider has no differential to expose; the shared common part stays in $\rng(\Gmat)$ and is
undetectable at any provider count (needs an external tie).}
\label{tab:provider}
\small
\setlength{\tabcolsep}{10pt}
\begin{tabular}{@{}ccl@{}}
\toprule
\textbf{Providers} & \textbf{Detectable cross-provider rows} & \textbf{Capability} \\
\midrule
$1$ & $0$  & no internal cross-check \\
$2$ & $48$ & detect disagreement, cannot identify culprit \\
$4$ & $72$ & detect and identify culprit (block-spark $>2f$) \\
\bottomrule
\end{tabular}
\end{table}
